\documentclass[10pt,a4paper]{article}
\usepackage{amsmath,amssymb,amsthm,graphicx}
\usepackage{amsthm}
\usepackage{amsfonts, amsmath, amssymb, amsthm}
\usepackage{hyperref}
\usepackage{cleveref}
\def\TT{\mathbb{T}}
\def\PP{\mathbb{P}}
\def\DD{\mathbb{D}}
\def\QQ{\mathbb{Q}}
\def\WW{\mathbb{W}}
\def\MM{\mathbb{M}}
\def\HH{\mathbb{H}}

\newcommand{\Psf}{\mathsf{P}}
\newcommand{\Qsf}{\mathsf{Q}}
\newcommand{\Rsf}{\mathsf{R}}


\def\F{\mathcal{F}}
\def\U{\mathcal{U}}
\def\C{\mathcal{C}}
\def\I{\mathcal{I}}
\def\D{\mathcal{D}}
\def\L{\mathcal{L}}
\def\V{\mathcal{V}}
\def\A{\mathcal{A}}
\def\B{\mathcal{B}}
\def\P{\mathcal{P}}
\def\M{\mathcal{M}}
\def\Nc{\mathcal{N}}
\def\Q{\mathcal{Q}}

% Baire space, Cantor space, binary strings
\newcommand{\cs}{2^{\NN}}
\newcommand{\bs}{\NN^{\NN}}
\newcommand{\bstr}{2^{<\NN}}

\newcommand{\uh}{\upharpoonright}

% Turing machines
\newcommand{\converge}{\!\!\downarrow}
\newcommand{\diverge}{\!\!\uparrow}

% Forcing question
\newcommand{\qvdash}{\operatorname{{?}{\vdash}}}
\newcommand{\nqvdash}{\operatorname{{?}{\nvdash}}}

\newcommand{\ISig}{\mathsf{I}\Sigma^0}
\newcommand{\Robinson}{\mathsf{Q}}
\newcommand{\IPi}{\mathsf{I}\Pi^0}
\newcommand{\IDelta}{\mathsf{I}\Delta^0}
\newcommand{\Iopen}{\mathsf{I_{open}}}
\newcommand{\BSig}{\mathsf{B}\Sigma^0}
\newcommand{\BPi}{\mathsf{B}\Pi^0}
\newcommand{\BDelta}{\mathsf{B}\Delta^0}
\newcommand{\BCSig}{\mathsf{BC}\Sigma^0}
\newcommand{\BCPi}{\mathsf{BC}\Pi^0}
\newcommand{\BCDelta}{\mathsf{BC}\Delta^0}
\newcommand{\LSig}{\mathsf{L}\Sigma^0}
\newcommand{\LPi}{\mathsf{L}\Pi^0}
\newcommand{\CSig}{\mathsf{C}\Sigma}

\newcommand{\RCA}[0]{\mathsf{RCA}}
\newcommand{\WKL}[0]{\mathsf{WKL}}
\newcommand{\WWKL}[0]{\mathsf{WWKL}}
\newcommand{\ACA}[0]{\mathsf{ACA}}
\newcommand{\ATR}[0]{\mathsf{ATR}}
\newcommand{\PIOOCA}[0]{\Pi^1_1\mathsf{-CA}}
\newcommand{\CAC}[0]{\mathsf{CAC}}
\newcommand{\ADS}[0]{\mathsf{ADS}}
\newcommand{\RRT}[0]{\mathsf{RRT}}
\newcommand{\FS}[0]{\mathsf{FS}}
\newcommand{\PFS}[0]{\mathsf{PFS}}
\newcommand{\TS}[0]{\mathsf{TS}}
\newcommand{\GRRT}[0]{\mathsf{GRRT}}
\newcommand{\LRRT}[0]{\mathsf{LRRT}}
\newcommand{\EM}[0]{\mathsf{EM}}
\newcommand{\SEM}[0]{\mathsf{SEM}}
\newcommand{\HalfEM}[0]{\mathsf{Half}\mbox{-}\mathsf{EM}}
\newcommand{\RT}[0]{\mathsf{RT}}
\newcommand{\SRT}[0]{\mathsf{SRT}}
%\newcommand{\RRT}{\mathsf{RRT}}

\newcommand{\NN}[0]{\mathbb{N}}
\newcommand{\nat}[0]{\mathbb{N}}

\newcommand{\card}{\operatorname{card}}
\newcommand{\finsub}{\subseteq_{\mathtt{fin}}}

% reductions

\newcommand{\W}{\mathrm{W}}
\newcommand{\sW}{\mathrm{sW}}

\newcommand{\lex}{\mathrm{lex}}

% Barriers
\newcommand{\base}{\operatorname{base}}
\newcommand{\range}{\operatorname{ran}}

%typographically correct double quotes
\def\qt#1{``#1''}%
\newtheorem*{statement}{Statement}
\newtheorem{theorem}{Theorem}
\numberwithin{theorem}{section}
\newtheorem{maintheorem}[theorem]{Main Theorem}
\newtheorem{lemma}[theorem]{Lemma}
\newtheorem{question}[theorem]{Question}
\newtheorem{proposition}[theorem]{Proposition}
%\newtheorem{statement}[theorem]{Statement}
\newtheorem{remark}[theorem]{Remark}
\newtheorem{definition}[theorem]{Definition}
\newtheorem{corollary}[theorem]{Corollary}
\newtheorem{example}[theorem]{Example}

\newcommand{\R}{\mathcal{R}}
\renewcommand{\S}{\mathcal{S}}
\newcommand{\T}{\mathcal{T}}
\title{A computable omega-to-omega barrier with progressive free sets that compute the halting set}
\author{Lorenzo Carlucci, Oriola Gjetaj, Quentin Le Houérou, Ludovic Levy Patey}
\date{arxiv:2412.11598v3; CC BY 4.0}
\begin{document}
\maketitle
\noindent\textit{Editorial note.} Solutions in the counted statement are infinite f-free sets, as in the progressive free-set principle1023-1026 and the proof1072. Finite free sets are not claimed to compute the jump.\par\medskip
\noindent\textit{Editorial note.} Normalize the modulus to be nondecreasing and at least3, which is harmless and ensures all finite thin-coloring instances lie in the proved n>=2 range.\par\medskip
\noindent\textit{Editorial note.} The finite-jump coloring array is shared support of the existing Schreier omega-jump task and is not counted again; the new object is a computable progressive coloring on a non-computably-bounded barrier.\par\medskip
\noindent\textit{Editorial scope.} One original result. Full source-local mathematical construction included; named classical inputs are explicit.\par\medskip
\noindent\textit{Allowed standard prerequisites (editorial).}\begin{itemize}
\item Classical computability tools. The halting set has a monotone left-c.e. modulus; the uniform relativized limit lemma removes one jump by one extra coordinate; a function dominating a modulus computes the set.
\item Classical finite-order barrier characterization. Assous PropositionII.1: a barrier has order type at most omega\textasciicircum{}n exactly when its member sizes are bounded by n. The source omega-bounded equivalence is supplied as local proof support.
\end{itemize}\medskip
There exist four other subsystems which, together with $\RCA_0$, calibrate the strength of most theorems. These subsystems are known as the \qt{Big Five} (see Mont\'alban~\cite{montalban_open_2011}). Among these, $\ACA_0$ (Arithmetic Comprehension Axiom) extends the axioms of $\RCA_0$ with the comprehension scheme for all arithmetic formulas. Based on the correspondence between computability and definability, it is equivalent to stating the existence of the Turing jump of any set ($\forall X \exists Y (Y = X')$).

We shall also consider two lesser-known stronger variants of $\ACA_0$, namely, $\ACA_0'$ and $\ACA_0^+$. The system $\ACA_0'$ extends $\RCA_0$ with the axiom stating closure under all finite jumps; the system $\ACA_0^+$ extends $\RCA_0$ with the axiom stating the existence of the $\omega$-jump of any set. The $\omega$-jump of a set~$X$ is the set $X^{(\omega)} = \bigoplus_n X^{(n)}$. The system $\ACA_0^+$ is famous for being the best known upper bound to Hindman's theorem.%\lorenzo{I added mention of ACA0' since we mention it at least once.}

Let us first fix some notation and recall the definition of these principles for colorings of finite subsets of fixed size. We denote by $\NN$ the set of natural numbers and by $\NN^+$ the set of positive natural numbers. For $X\subseteq \NN$ and $n \in \NN$
we denote by $[X]^{n}$ the set of all $n$-subsets of $X$. We denote by $[X]^{<\omega}$ the set of all finite subsets of $X$. We always assume that sets are presented in increasing order and make no distinction between a set and the sequence of its elements in increasing order. We write $x \cdot y$ for the concatenation of two sequences $x$ and $y$. We identify a natural number with the set of its predecessors, so that, if $k\geq 1$ we can write $f: X \to k$ to declare a function from $X$ to $\{0, 1, \dots, k-1\}$. A function of this type is often called a coloring of $X$ in $k$ colors.

\begin{definition}
Let $C$ be any non-empty set. Given a coloring $f : [\NN]^n \to C$, a set~$H \subseteq \NN$ is \emph{$f$-thin for color~$c \in C$} if $c \not \in f([H]^n)$. A set ~$H \subseteq \NN$ is $f$-thin (or thin for $f$) if $H$ is $f$-thin for some color $c\in C$.
\end{definition}

The definition of $f$-thin set depends on the choice of codomain of the function. It will always be clear from the context. 

\begin{definition}
A set $B \subseteq [\NN]^{<\omega}$ is a \emph{barrier} if
\begin{enumerate}
    \item[(1)] $\base(B)$ is infinite;
    \item[(2)] for every~$X \in [\base(B)]^\omega$, there is some~$s \in B$ such that $s \preceq X$;
    \item[(3)] for every~$s, t \in B$, $s \not \subset t$.
\end{enumerate}
\end{definition}

Given a set~$X$ and $n \in \NN$, we write $[X]^{\leq n}$ for the set of all subsets $s \subseteq X$ such that $|s| \leq n$, and $[X]^{\leq !\omega}$ for the set of all subsets $s \subseteq X$ such that $|s| \leq 1+\min s$. By convention, for~$n = 0$, $[X]^{\leq n}$ is the singleton $\{\emptyset\}$. Given a function $h : \NN \to \NN$, we write $[X]^{\leq h(\cdot)}$ for the set of all finite $s \subseteq X$ such that $|s| \leq h(\min s)$. %\lorenzo{why don't we just write $[X]^{\leq n}$ for $[X]^{!\leq n}$? Same for $\leq h$}\ludovic{This is legacy from when I was using $\alpha \leq \omega$. we can remove the exclamation point. I will do it if you want.}
%\lorenzo{I find notation $[X]^{\leq h(\cdot)}$ a bit confusing. Is $[X]^{\leq h(\cdot)}$ viable for you?}\ludovic{Yes, I will change it}

\begin{definition}
A set $B \subseteq [\NN]^{<\omega}$ is \emph{$\omega$-bounded} if and only if $B \subseteq [\NN]^{\leq h(\cdot)}$ for some function~$h : \NN \to \NN$. It is \emph{computably $\omega$-bounded} if furthermore $h$ is computable.
\end{definition}

A function $h : \NN \to \NN$ is \emph{left-c.e.} if there is a uniformly computable sequence of functions $h_0, h_1, \dots$ such that for every~$x, i \in \NN$, $h_i(x) \leq h_{i+1}(x)$ and $\lim_i h_i(x) = h(x)$. The sequence $(h_i)_{i\in\NN}$ is then called a \emph{left-c.e. approximation} of $h$.
If $B$ is a computable barrier of order type at most $\omega^\omega$, then it is $\omega$-bounded by a left-c.e. function. This bound is tight, as there exist computable barriers of order type $\omega^{\omega}$ which are not computably $\omega$-bounded.

\begin{definition}
Let $B \subseteq [\NN]^{<\omega}$ be a set and $f : B \to [\NN]^{<\omega}$ be a coloring.
A set~$H \subseteq \base(B)$ is \emph{$f$-free} if for every~$s \in B \cap [H]^{<\omega}$, $f(s) \cap H \subseteq s$.
\end{definition}

\begin{definition}
Fix $B \subseteq [\NN]^{<\omega}$.
A function $f : B \to [\NN]^{<\omega}$ is \emph{progressive} if for every~$s \in B$, either $f(s) = \emptyset$, or $\min f(s)\geq \min s$.
\end{definition}

An easy combinatorial argument shows that the following statement is classically true.


\begin{definition}[Progressive Free Set Theorem]
Fix a set $B \subseteq [\NN]^{<\omega}$ and a bounding function $b : \NN \to \NN$.
$\PFS^B_b$ is the statement \qt{For every $b$-constrained progressive coloring $f : B \to [\NN]^{<\omega}$, there is an infinite $f$-free set}.
\end{definition}

\begin{proposition}\label[proposition]{prop:computable-barrier-pfs-computes-zp}
    There exists a computable barrier $B$ of order type $\omega^{\omega}$ and a computable progressive coloring $f : B \to \NN$ such that every $f$-free set computes~$\emptyset'$.
\end{proposition}

\begin{lemma}\label{lem:arrays}
    There exists two computable arrays $(e_{n,k})_{n,k \in \NN}$ and $(d_{n,k})_{n,k \in \NN}$ of Turing indexes such that for every $n, k \in \NN^+$ with $n \geq 2$, $\Phi_{e_{n,k}}^{\emptyset^{(k)}}$ is a coloring $f_n^k : [\NN]^{n} \to n$ such that for every infinite $f_n^k$-thin set $H$, $\Phi_{d_{n,k}}^{H} = \emptyset^{(k)}$.
\end{lemma}

\begin{proof}
Let $k\in\NN^+$. Let $g_k$ be a uniform modulus of the set $\emptyset^{(k)}$. Note that $\Delta^0_{k+1}$-indexes for each function $g_k$ can be found computably, uniformly in~$k$. 

For $n \geq 2$ let $f_n^k(x_0, \dots,x_{n-1}) = n - 1$ if $g_k(x_0) \leq x_1$ and if $g_k(x_0) > x_1$, let $f_n^k(x_0, \dots,x_{n-1}) = i$ for $i < n - 1$ the largest value such that $g_k(x_0) > x_{i+1}$.

Let $H$ be an infinite $f_n^k$-thin set for some color $c$. The color $c$ cannot be $n-1$ as for a given $x_0 \in H$ there exists $x_1 <\dots < x_{n-1} \in H \setminus \{0, \dots, g_k(x_0)\}$, hence $f^k_n(x_0, \dots, x_{n-1}) = n-1$.

We claim $H$ is thin for color $n-2$. Indeed, assume by contradiction that there exists some tuple $x_0 < \dots < x_{n-1} \in H$ such that $g_k(x_0) > x_{n-1}$ so that $f^k_n$ takes color $n-2$ on $\{x_0, \dots, x_{n-1}\}$. Since $c < n-2$ and $H$ is infinite, there exists some $y_{c+2}< \dots < y_{n-1} \in H \setminus \{0, \dots ,g_k(x_0)\}$, therefore $f^k_n(x_0, \dots, x_{c+1},y_{c+2}, \dots, y_{n-1}) = c$, contradicting our assumption that $H$ is $f_n^k$-thin for $c$.

Write $H = \{x_0 < x_1 < \dots \}$. For every $i \in \NN$, $g_k(x_i) \leq x_{i + n - 1}$, hence $H$ computes a function dominating $g_k$ and therefore computes $\emptyset^{(k)}$. This computation can be done uniformly in the set $H$, $k$ and $n$. 
\end{proof}

\begin{lemma}
    There exists a functional $\Gamma$ such that for every set $X$ and every index $e$, if $\Phi_e^{X'}$ is a coloring $f : [\NN]^n \to \ell$, then $\Gamma^X(e)$ is a coloring $g : [\NN]^{n+1} \to \ell$ such that if an infinite set $H$ is $g$-thin for a color $c$, then $H$ is $f$-thin for the same color $c$. 
\end{lemma}

\begin{proof}
Consider a set $X$ and an index $e$ such that $\Phi_e^{X'}$ is a coloring $f : [\NN]^n \to \ell$ for some $n, \ell \in \NN^+$. Let $(f_s)_{s\in \NN}$ be a $\Delta_2^0(X)$-approximation of $f$ (indexes for such an approximation, where the $f_s$ are seen as $X$-computable functions, can computably be found uniformly in $X$ and $e$). Finally, consider the following $X$-computable coloring $g$ defined by $g(x_0, \dots, x_{n}) = f_{x_n}(x_0, \dots, x_{n-1})$ (again, the construction is uniform).

Let $H \subseteq \NN$ be an infinite set. If for some $x_0 < \dots < x_{n-1} \in H$ we have $f(x_0, \dots, x_{n-1}) = c$ for some color $c$, then $f_s(x_0, \dots, x_{n-1}) = c$ for every $s$ bigger than a certain threshold. Thus, as $H$ is infinite, there exists some $x_n \in H$ such that $g(x_0, \dots, x_n) = c$.  
\end{proof}

Combining these two lemmas, we get the following, where the array $(d_{n,k})_{n,k\in\NN}$ is as in~\Cref{lem:arrays}.

\begin{lemma}\label[lemma]{lem:uniform-computable-thin-coloring}
    There exists a uniformly computable sequence of colorings $f_{n,k} : [\NN]^{n+k} \to n$, for $n, k \in \NN^+$, such that for every infinite $f_{n,k}$-thin set $H$, $\Phi_{d_{n,k}}^H = \emptyset^{(k)}$.
\end{lemma}

\begin{lemma}\label[lemma]{barrier-omega-bounded-order type}
A barrier~$B$ has order type at most $\omega^{\omega}$ if and only if $B$ is $\omega$-bounded.
\end{lemma}
\begin{proof}
Suppose first $B \subseteq [\NN]^{\leq h(\cdot)}$ for some function $h : \NN \to \NN$.
For every~$s \in B$, let $\alpha_s = \sum_{i < |s|} \omega^{h(\min s)-i} s(i)$. %\lorenzo{$h(\min s)-i$ can be negative.}\ludovic{No, it cannot, since $B \subseteq [\NN]^{\leq h(\cdot)}$, so $|s| \leq h(\min s)$.}\lorenzo{I realized that - sorry ;)}
Note that $s <_{\lex} t$ if and only if $\alpha_s < \alpha_t$, so $B$ has order type at most $\omega^\omega$ since $\alpha_s < \omega^{\omega}$ for every $s \in B$.

Suppose now $B$ has order type at most $\omega^{\omega}$. Let $x \in \NN$ and $B_x = \{ s : x \cdot s \in B \}$. Then $B_x$ is a barrier of order type at most $\omega^{n_x}$ for some~$n_x \in \NN^+$. By Assous~\cite[Proposition II.1]{assous1974caracterisation}, a barrier~$B$ has order type at most $\omega^n$ if and only if $B \subseteq [\NN]^{\leq n}$, so $B_x \subseteq [\NN]^{\leq n_x}$. Let $h(x) = n_x$. Then $B \subseteq [\NN]^{\leq h(\cdot)}$.
\end{proof}

\begin{proof}
Let $g : \NN \to \NN$ be the modulus of $\emptyset'$ and let $(g_n)_{n \in \NN}$ be a left-c.e. approximation of $g$. Without loss of generality, it can be assumed that $g_{n}(x) \geq 1$ for every $x,n \in \NN$ and that $g_n$ is non-decreasing for each~$n \in \NN$.

Let $B$ be defined as follows: for $x \in \NN$ and $s \subseteq \NN$ with $x < \min s$, let $x \cdot s \in B$ if and only if $|s| = g_{\min s}(x)$. 

%\lorenzo{small problem here if $g_{\min s}(x)=0$. Can we just use $\NN^+$ as codomain in the definition of left-c.e. approximation or just assume here wlog that every value is $\geq 1$?}\ludovic{Quentin, I let you handle this. Let's rather assume wlog that every value is $\geq 1$.}\lorenzo{done} 

We claim that $B$ is a barrier with base $\NN$: for any infinite subset $X = \{x_0, x_1, \dots\}$ of $\NN$, we have $\{x_0, x_1, \dots, x_{g_{x_1}(x_0)}\}\in B$ (since $g_{x_1}(x_0)\geq 1$). Let $x \cdot s, y \cdot t$ be in $B$ and suppose $x\cdot s \subseteq y \cdot t$.
Then $x \geq y$ and $\min s \geq \min t$. Moreover $|s| = g_{\min s}(x) \geq g_{\min s}(y)\geq g_{\min t}(y) = |t|$. Thus $s = t$ and $x=y$. By \Cref{barrier-omega-bounded-order type}, the order type of~$B$ is $\omega^{\omega}$.

Let $f : B \to \NN$ be defined by $f(x \cdot s) = x + h_{|s| - 1}(s) + 1$ where $h_n : [\NN]^{n+1} \to n$ is the computable instance of $\TS^{n+1}_n$ obtained in \Cref{lem:uniform-computable-thin-coloring} such that every infinite $h_n$-thin set computes $\emptyset'$. 
Let $H$ be an infinite $f$-free set. For every $x \in H$, there exists some bound $b_x > g(x)$ such that $g_{b_x}(x) = g(x)$. For every $s \subseteq H \setminus [0,b_x]$ of cardinality $g(x)$, we have $f(x \cdot s) = x + h_{g(x) - 1}(s) + 1$. There are two cases:
\smallskip

\textbf{Case 1:} $H \setminus [0,b_x]$ is $h_{g(x) - 1}$-thin for some $x \in H$, in that case, by definition of $h_{g(x) - 1}$, $H \geq_T \emptyset'$.
\smallskip

\textbf{Case 2:} $H \setminus [0,b_x]$ is not $h_{g(x)}$-thin for every $x \in H$, in that case, for every $x \in H$ and every $y \in (x,x + g(x))$, there exists some $s \subseteq H \setminus [0,b_x]$ such that $f(x \cdot s) = y$ and therefore, since $H$ is $f$-free, $y$ is not in $H$. So $H \geq_T \emptyset'$ as the principal function of $H$ dominates $g$.
\end{proof}

\bibliographystyle{plain}
\bibliography{biblio}
\end{document}
